\section{Problem Formulation}

\textbf{Motivation.} Fashion-MNIST is a standard image-classification benchmark for evaluating machine-learning pipelines on grayscale clothing images. The task provides a controlled setting for comparing a linear classifier with a multilayer perceptron while preserving challenges such as visually similar clothing categories.

\textbf{Input.} Each input is a grayscale image with shape $1\times28\times28$. Pixel intensities are represented as floating-point values in $[0,1]$ before normalization and are flattened to 784 features by both models.

\textbf{Output.} The target is one of ten Fashion-MNIST classes. Each prediction corresponds to one image and is obtained from the class with the largest output logit.

\textbf{Mathematical formulation.}
\begin{equation}
    \hat{y} = f_{\theta}(x),
    \qquad
    \theta^{*} = \arg\min_{\theta}\frac{1}{N}\sum_{i=1}^{N}
    \mathcal{L}\bigl(f_{\theta}(x_i), y_i\bigr).
    \label{eq:problem-formulation}
\end{equation}
For a batch of images, the models minimize the cross-entropy loss over the ten-class target distribution.

\textbf{Main difficulties.} The main challenges are distinguishing visually similar classes such as Shirt, T-shirt/top, Pullover, and Coat; handling the high-dimensional flattened input; and controlling overfitting while using a fully connected architecture.
